\documentclass{article}

\usepackage[margin=1in,letterpaper]{geometry}
\usepackage{amsfonts,amsmath,amsthm}
\usepackage{xcolor}
\usepackage{algpseudocodex}
\usepackage[package=algpseudocodex]{texfrog}

% Minimal macros for the TeXFrog IND-CPA tutorial example (algpseudocodex version).
% These supplement the commands provided by the algpseudocodex package.

\newcommand{\INDCPA}{\mathrm{IND\text{-}CPA}}
\newcommand{\Enc}{\mathsf{Enc}}
\newcommand{\Adversary}{\mathcal{A}}
\newcommand{\Bdversary}{\mathcal{B}}
\newcommand{\OPRF}{\mathcal{O}_{\mathsf{prf}}}
\newcommand{\RF}{\mathrm{RF}}
\newcommand{\getsr}{{\:\leftarrow\hspace*{-3pt}\raisebox{.5pt}{$\scriptscriptstyle\$$}\:}}


\newtheorem{theorem}{Theorem}

%%% Game registration
\tfgames{indcpa}{G0, G1, Red1, G2, G3}
\tfgamename{indcpa}{G0}{G_0}
\tfgamename{indcpa}{G1}{G_1}
\tfgamename{indcpa}{Red1}{\Bdversary_1}
\tfgamename{indcpa}{G2}{G_2}
\tfgamename{indcpa}{G3}{G_3}

\tfdescription{indcpa}{G0}{$\INDCPA_\Enc^\Adversary()$: The IND-CPA game. The LR oracle encrypts using the actual PRF.}
\tfdescription{indcpa}{G1}{Replace $\mathrm{PRF}(k, r)$ with a truly random function $\RF(r)$.}
\tfdescription{indcpa}{Red1}{Reduction against PRF security, simulating Games \tfgamename{indcpa}{G0} and \tfgamename{indcpa}{G1} via an external PRF challenger.}
\tfdescription{indcpa}{G2}{Replace $\RF(r)$ with independently sampled randomness.}
\tfdescription{indcpa}{G3}{Ciphertext is uniformly random, independent of $m_b$.}

\tfreduction{indcpa}{Red1}
\tfrelatedgames{indcpa}{Red1}{G0, G1}

\tfmacrofile{macros.tex}

\tfcommentary{indcpa}{G0}{commentary/G0.tex}
\tfcommentary{indcpa}{G1}{commentary/G1.tex}
\tfcommentary{indcpa}{Red1}{commentary/Red1.tex}
\tfcommentary{indcpa}{G2}{commentary/G2.tex}
\tfcommentary{indcpa}{G3}{commentary/G3.tex}

\tffigure{indcpa}[Games $\tfgamename{indcpa}{G0}$--$\tfgamename{indcpa}{G3}$, Reduction $\tfgamename{indcpa}{Red1}$]{all_games}{G0,G1,G2,G3,Red1}

%%% Proof source
\begin{tfsource}{indcpa}
\begin{algorithmic}[1]
\Procedure{%
  \tfonly*{G0}{$\INDCPA_\Enc^\Adversary()$}%
  \tfonly*{G1}{Game~$\tfgamename{G1}$}%
  \tfonly*{G2}{Game~$\tfgamename{G2}$}%
  \tfonly*{G3}{Game~$\tfgamename{G3}$}%
  \tfonly*{Red1}{Reduction $\Bdversary_1^{\OPRF}$}%
  \tffigonly{Games $\tfgamename{G0}$--$\tfgamename{G3}$}%
}{}
\tfonly{G0}{\State $k \getsr \{0,1\}^\lambda$}
\State $b \getsr \{0,1\}$
\State $b' \gets \Adversary^{\mathsf{LR}}()$
\State \Return $(b' = b)$
\EndProcedure
\Statex
\Procedure{Oracle $\mathsf{LR}(m_0, m_1)$}{}
\State $r \getsr \{0,1\}^\lambda$
\tfonly{G0}{\State $y \gets \mathrm{PRF}(k, r)$}
\tfonly{G1}{\State $y \gets \RF(r)$}
\tfonly{G2}{\State $y \getsr \{0,1\}^\lambda$}
\tfonly{Red1}{\State $y \gets \OPRF(r)$}
\tfonly{G0-G2}{\State $c \gets y \oplus m_b$}
\tfonly{G3}{\State $c \getsr \{0,1\}^\lambda$}
\State \Return $(r, c)$
\EndProcedure
\end{algorithmic}
\end{tfsource}

\title{IND-CPA Security of PRF-Based Symmetric Encryption\\
{\large A TeXFrog Tutorial (using the \texttt{algpseudocodex} package)}}
\author{}
\date{}

\begin{document}

\maketitle

\section{Construction}

Let $\mathrm{PRF}\colon \{0,1\}^\lambda \times \{0,1\}^\lambda \to \{0,1\}^\lambda$
be a pseudorandom function.
Define a symmetric encryption scheme $\Pi = (\mathsf{KeyGen}, \mathsf{Enc}, \mathsf{Dec})$ as follows:
\begin{itemize}
  \item $\mathsf{KeyGen}()$: sample $k \getsr \{0,1\}^\lambda$ and return $k$.
  \item $\mathsf{Enc}(k, m)$: sample $r \getsr \{0,1\}^\lambda$, compute
    $c \gets \mathrm{PRF}(k, r) \oplus m$, and return $(r, c)$.
  \item $\mathsf{Dec}(k, (r, c))$: return $\mathrm{PRF}(k, r) \oplus c$.
\end{itemize}

\section{Security}

\begin{theorem}
If\/ $\mathrm{PRF}$ is a secure pseudorandom function, then $\Pi$ is IND-CPA secure.
Concretely, for any adversary $\mathcal{A}$ making $q$ queries to the
left-or-right oracle,
\[
  \mathbf{Adv}^{\mathrm{IND\text{-}CPA}}_\Pi(\mathcal{A})
  \;\leq\;
  \mathbf{Adv}^{\mathrm{PRF}}_{\mathrm{PRF}}(\mathcal{B}_1)
  \;+\; \frac{q^2}{2^{\lambda+1}} \,.
\]
\end{theorem}

\noindent\textit{Proof.}\quad
We proceed via a sequence of games. Figure~\ref{fig:all-games} shows all the games, and the games are also shown individually in subsequent figures.

\begin{figure}[ht]
\centering
\tfrenderfigure{indcpa}{G0,G1,G2,G3,Red1}
\caption{Consolidated figure showing games $\tfgamename{indcpa}{G0}$--$\tfgamename{indcpa}{G3}$ and Reduction $\tfgamename{indcpa}{Red1}$.}
\label{fig:all-games}
\end{figure}

\paragraph{Game~$\tfgamename{indcpa}{G0}$.}
The starting game is $\tfgamename{indcpa}{G0} = \INDCPA_\Enc^\Adversary()$.
The challenger samples a uniform key $k$ and challenge bit $b$, then
answers each $\mathsf{LR}(m_0, m_1)$ query with a fresh encryption
$(r, \mathrm{PRF}(k,r) \oplus m_b)$.


\begin{figure}[ht]
\centering
\tfrendergame{indcpa}{G0}
\caption{Game~$\tfgamename{indcpa}{G0}$: $\INDCPA_\Enc^\Adversary()$: The IND-CPA game. The LR oracle encrypts using the actual PRF.}
\label{fig:G0}
\end{figure}

\paragraph{Game~$\tfgamename{indcpa}{G1}$.}
\textbf{Claim.}
  Games \tfgamename{indcpa}{G0} and \tfgamename{indcpa}{G1} are computationally indistinguishable assuming $\mathrm{PRF}$
  is a secure pseudorandom function, with advantage gap at most
  $\mathbf{Adv}^{\mathrm{PRF}}(\Bdversary_1)$.

Reduction \tfgamename{indcpa}{Red1} (shown separately) interacts with an external PRF challenger
and simulates the LR oracle by forwarding nonces $r$ to that challenger.
When the external challenger uses the real PRF, \tfgamename{indcpa}{Red1} exactly simulates
\tfgamename{indcpa}{G0}; when it uses a random function, \tfgamename{indcpa}{Red1} exactly simulates \tfgamename{indcpa}{G1}.


\begin{figure}[ht]
\centering
\tfrendergame[diff=G0]{indcpa}{G1}
\caption{Game~$\tfgamename{indcpa}{G1}$: Replace $\mathrm{PRF}(k, r)$ with a truly random function $\RF(r)$.}
\label{fig:G1}
\end{figure}

\paragraph{Reduction~$\tfgamename{indcpa}{Red1}$.}
The reduction \tfgamename{indcpa}{Red1} receives no PRF key; instead it routes each nonce $r$
to the external PRF oracle $\OPRF$.  All other bookkeeping (sampling $b$,
running $\Adversary$, checking the guess) is identical to Games \tfgamename{indcpa}{G0} and \tfgamename{indcpa}{G1}.


\begin{figure}[ht]
\centering
\tfrendergame[diff=G1]{indcpa}{Red1}
\caption{Reduction~$\tfgamename{indcpa}{Red1}$: Reduction against PRF security, simulating Games~\tfgamename{indcpa}{G0} and~\tfgamename{indcpa}{G1} via an external PRF challenger.}
\label{fig:Red1}
\end{figure}

\paragraph{Game~$\tfgamename{indcpa}{G2}$.}
\textbf{Claim.}
  Games \tfgamename{indcpa}{G1} and \tfgamename{indcpa}{G2} are statistically close:
  $|\Pr[\tfgamename{indcpa}{G1} = 1] - \Pr[\tfgamename{indcpa}{G2} = 1]| \leq q^2 / 2^{\lambda+1}$.

In \tfgamename{indcpa}{G1}, each query samples a fresh nonce $r \getsr \{0,1\}^\lambda$
and evaluates the random function $\RF$ at $r$.  If all $q$ nonces are distinct---which
fails with probability at most $\binom{q}{2}/2^\lambda = q(q-1)/2^{\lambda+1}$
by the birthday bound---then the outputs $\RF(r_1), \ldots, \RF(r_q)$
are independent and uniformly distributed, exactly as in \tfgamename{indcpa}{G2}.


\begin{figure}[ht]
\centering
\tfrendergame[diff=G1]{indcpa}{G2}
\caption{Game~$\tfgamename{indcpa}{G2}$: Replace $\RF(r)$ with independently sampled randomness.}
\label{fig:G2}
\end{figure}

\paragraph{Game~$\tfgamename{indcpa}{G3}$.}
\textbf{Claim.}
  Games \tfgamename{indcpa}{G2} and \tfgamename{indcpa}{G3} are perfectly indistinguishable.

In \tfgamename{indcpa}{G2}, each query samples a uniform $y$ independently.
Since $c = y \oplus m_b$ with uniform $y$, the ciphertext $c$ is
uniformly distributed regardless of $m_b$.
Game \tfgamename{indcpa}{G3} makes this explicit by sampling $c$ directly.

In Game \tfgamename{indcpa}{G3} the adversary's view is identically distributed for $b = 0$ and $b = 1$,
so $\Pr[b' = b] = \tfrac{1}{2}$ and the IND-CPA advantage is~$0$.


\begin{figure}[ht]
\centering
\tfrendergame[diff=G2]{indcpa}{G3}
\caption{Game~$\tfgamename{indcpa}{G3}$: Ciphertext is uniformly random, independent of~$m_b$.}
\label{fig:G3}
\end{figure}

\noindent
Combining the bounds from each game hop, we obtain
\[
  \mathbf{Adv}^{\mathrm{IND\text{-}CPA}}_\Pi(\mathcal{A})
  \;=\;
  \bigl|\Pr[\mathsf{G}_0 = 1] - \tfrac{1}{2}\bigr|
  \;\leq\;
  \mathbf{Adv}^{\mathrm{PRF}}_{\mathrm{PRF}}(\mathcal{B}_1)
  \;+\; \frac{q^2}{2^{\lambda+1}}
  \;+\; 0 \,,
\]
which completes the proof.\hfill$\square$

\end{document}
